\documentclass[11pt,letterpaper]{letter}
\usepackage[utf8]{inputenc}
\usepackage[T1]{fontenc}
\usepackage[english]{babel}
\usepackage{geometry}
\geometry{left=1in,right=1in,top=1in,bottom=1in}
\usepackage{setspace}
\onehalfspacing
\usepackage[colorlinks=true,urlcolor=blue,linkcolor=black]{hyperref}

\name{Ángel Luis Robles Fernández}
\address{5125 Overland Dr. B8, 66049\\
Lawrence, KS\\
+1 480 791 4846\\
a.l.robles.fernandez@gmail.com\\
\url{https://alrobles.github.io/}}
\signature{Ángel Luis Robles Fernández}
\date{October 1, 2026}

\begin{document}
\begin{letter}{
Search Committee for Assistant Professor\\
Department of Integrative Biology\\
Oregon State University\\
2403 Cordley Hall\\
Corvallis, OR 97331
}

\opening{Dear Members of the Integrative Biology Search Committee,}

I am applying for the tenure-track Assistant Professor search in Integrative Biology (JR0000039). I earned my Ph.D. in Evolutionary Biology at Arizona State University in 2025 and am a postdoctoral researcher in Daniel Reuman's laboratory at the University of Kansas. I would contribute a quantitative research program connecting environmental tolerance and organismal performance information with population persistence, supported by mathematical modeling, machine learning, and open research software.

My work in \emph{PNAS} developed frameworks for predicting wildlife host--pathogen associations across contrasting disease systems. As co-first author of an \emph{Ecography} study, I linked historical environmental suitability with population genetic diversity. At Kansas, I co-developed XSDM, which incorporates stochastic demography and climate variability into distribution modeling, and created and maintain its R package. I also develop its Bayesian implementation, \texttt{xsdmbayes}. These contributions demonstrate my ability to turn biological questions into quantitative frameworks and reusable scientific products.

An ongoing collaboration provides a concrete bridge to the department's emphasis on organismal biology. We propose using experimental performance--environment curves to inform Bayesian XSDM priors and comparing models with and without that experimental information. This work has not yet been published. At OSU, I will build on this direction to ask how environmental-performance relationships and climate variability jointly determine persistence. Working alongside colleagues who characterize organismal responses, I will lead statistical integration and software development, testing when experimental information improves geographic predictions or exposes limits to their transferability.

I will also pursue AI-enabled biological databases and rigorously evaluated AI-driven scientific workflows. My stack includes \href{https://ecoseek.org}{EcoSeek}, \href{https://genominer.ecoseek.org}{Genominer}, and \href{https://xbioclim.org}{xbioclim}, connecting ecological analyses, population-genetic resources, and time-resolved climate data. DiDAL, litreview, and Sciev extend this agenda toward scientist-led synthesis and structured evidence assessment. My background in physics, evolutionary biology, and industrial machine learning positions me to coordinate interdisciplinary research while pursuing external funding for integrative modeling, data infrastructure, and computational biological methods.

My teaching experience includes Introduction to R and teaching assistantships in Quantitative Ecology and Statistical Models for Biology at Arizona State. I have also been invited to teach graduate-level ecological niche modeling at INECOL and am an RStudio Certified Trainer in the Tidyverse. I use live coding, scaffolded projects, and collaborative debugging to connect theory with evidence. At OSU, I will combine active learning and formative assessment with accessible resources and support for varied preparation. Students using AI will disclose assistance, verify evidence, and explain their reasoning.

I am prepared to contribute to introductory biology, ecology, evolution, and quantitative courses in the biology and zoology programs. I will mentor undergraduate and graduate researchers through clear expectations, supported entry projects, and increasing scientific independence. My industry experience will also help students explore careers beyond academia. I would welcome the opportunity to build a collaborative research group that connects organismal expertise with ecological and evolutionary inference.

My application includes my CV, research statement, and teaching statement. Thank you for considering my application.

\closing{Sincerely,}
\end{letter}
\end{document}
